\documentclass{article}
\usepackage{amsmath,amssymb,amsthm}
\usepackage{algorithm}
\usepackage[noend]{algpseudocode}
\usepackage{fullpage}
\usepackage{hyperref}
\newtheorem{lemma}{Lemma}
\newtheorem{theorem}{Theorem}
\newtheorem{assumption}{Assumption}
\renewcommand{\thealgorithm}{}
\begin{document}

%=============================================================
\section*{Algorithm Name}
\textbf{Stochastic Variance-Reduced Gradient (SVRG) under the PL condition}
%=============================================================

%=============================================================
\section*{Mathematical Formulation}
Let $f(x)=\frac{1}{n}\sum_{i=1}^{n}f_i(x)$ be a finite‑sum objective, where each $f_i:\mathbb{R}^d\to\mathbb{R}$ is differentiable.
The algorithm proceeds in \emph{epochs}.  
For epoch $s=0,1,2,\dots$ we denote the \emph{snapshot point} by $\tilde{x}^{\,s}$ and its full gradient 
\[
\tilde{\mu}^{\,s}= \nabla f(\tilde{x}^{\,s}) = \frac{1}{n}\sum_{i=1}^{n}\nabla f_i(\tilde{x}^{\,s}).
\]

Inside epoch $s$ we perform $m$ inner stochastic updates.  
For $t=0,1,\dots,m-1$ let $x_t^{\,s}$ be the inner iterate (with $x_0^{\,s}=\tilde{x}^{\,s}$) and draw an index $i_t$ uniformly from $\{1,\dots,n\}$.  
Define the variance‑reduced estimator
\[
v_t^{\,s}= \nabla f_{i_t}(x_t^{\,s})-\nabla f_{i_t}(\tilde{x}^{\,s})+\tilde{\mu}^{\,s}.
\]
The inner update is a standard SGD step
\[
x_{t+1}^{\,s}= x_t^{\,s}-\eta\, v_t^{\,s},
\]
where $\eta>0$ is a constant stepsize.  
After $m$ inner steps we set the next snapshot
\[
\tilde{x}^{\,s+1}=x_m^{\,s}.
\]

The complete algorithm is summarised in the pseudocode below.  

%=============================================================
\section*{Pseudocode}
\begin{algorithm}[H]
\caption{SVRG under the PL condition}
\begin{algorithmic}[1]
\Require Number of epochs $S$, inner loop length $m$, stepsize $\eta>0$, initial point $x_0$.
\State $\tilde{x}^{\,0}\gets x_0$
\For{$s=0$ \textbf{to} $S-1$}
    \State Compute full gradient $\displaystyle \tilde{\mu}^{\,s}\gets \frac{1}{n}\sum_{i=1}^{n}\nabla f_i(\tilde{x}^{\,s})$
    \State $x_0^{\,s}\gets \tilde{x}^{\,s}$
    \For{$t=0$ \textbf{to} $m-1$}
        \State Sample $i_t\sim\text{Uniform}\{1,\dots,n\}$
        \State $v_t^{\,s}\gets \nabla f_{i_t}(x_t^{\,s})-\nabla f_{i_t}(\tilde{x}^{\,s})+\tilde{\mu}^{\,s}$
        \State $x_{t+1}^{\,s}\gets x_t^{\,s}-\eta\,v_t^{\,s}$
    \EndFor
    \State $\tilde{x}^{\,s+1}\gets x_m^{\,s}$
\EndFor
\State \Return $\tilde{x}^{\,S}$
\end{algorithmic}
\end{algorithm}

%=============================================================
\section*{Dry Run Example}
Consider the scalar quadratic 
\[
f(w)=(w-3)^2,\qquad n=1,\; f_1(w)=f(w).
\]
The gradient is $\nabla f(w)=2(w-3)$, which is $L=2$ smooth and satisfies the PL condition with $\mu=2$.

Take $x_0=0$, stepsize $\eta=0.1$, inner length $m=3$, and run a single epoch ($S=1$).

\begin{center}
\begin{tabular}{c|c|c|c}
$t$ & $x_t$ & $v_t$ (since $n=1$, $v_t=\nabla f(x_t)$) & $f(x_t)$\\ \hline
0 & $0$ & $2(0-3)=-6$ & $9$\\
1 & $0-0.1(-6)=0.6$ & $2(0.6-3)=-4.8$ & $(0.6-3)^2=5.76$\\
2 & $0.6-0.1(-4.8)=1.08$ & $2(1.08-3)=-3.84$ & $(1.08-3)^2=3.6864$\\
3 & $1.08-0.1(-3.84)=1.464$ & $2(1.464-3)=-3.072$ & $(1.464-3)^2=2.371$\\
\end{tabular}
\end{center}

After the inner loop we set $\tilde{x}^{\,1}=x_3=1.464$, and the objective has decreased from $9$ to $\approx2.37$ in three stochastic steps.

%=============================================================
\section*{Assumptions}
\begin{assumption}[Smoothness]\label{ass:smooth}
Each component $f_i$ is $L$‑smooth, i.e.
\[
\|\nabla f_i(x)-\nabla f_i(y)\|\le L\|x-y\|,\qquad\forall x,y\in\mathbb{R}^d .
\]
Consequently $f$ is $L$‑smooth as well.
\end{assumption}

\begin{assumption}[Polyak–Łojasiewicz (PL) condition]\label{ass:pl}
The objective $f$ satisfies
\[
\frac{1}{2}\|\nabla f(x)\|^2\ge \mu\bigl(f(x)-f^\star\bigr),\qquad\forall x\in\mathbb{R}^d,
\]
where $f^\star=\min_x f(x)$ and $\mu>0$.
\end{assumption}

\begin{assumption}[Stepsize]\label{ass:stepsize}
The stepsize $\eta$ is chosen such that
\[
0<\eta\le \frac{1}{3L}.
\]
\end{assumption}

\begin{assumption}[Inner loop length]\label{ass:inner}
The number of inner iterations $m$ satisfies
\[
m\ge \frac{2}{\eta\mu}.
\]
\end{assumption}

%=============================================================
\section*{Main Convergence Theorem}
\begin{theorem}[Linear convergence of SVRG under PL]\label{thm:linear}
Assume \ref{ass:smooth}–\ref{ass:inner} hold. Let $\tilde{x}^{\,s}$ be the snapshot after epoch $s$.
Define $\Delta_s:=\mathbb{E}\!\bigl[f(\tilde{x}^{\,s})-f^\star\bigr]$.
Then for every epoch $s\ge0$,
\[
\Delta_{s+1}\;\le\;\bigl(1-\tfrac{\eta\mu}{2}\bigr)\,\Delta_s .
\]
Consequently, after $S$ epochs,
\[
\mathbb{E}\!\bigl[f(\tilde{x}^{\,S})-f^\star\bigr]\;\le\;
\bigl(1-\tfrac{\eta\mu}{2}\bigr)^{S}\,\bigl(f(x_0)-f^\star\bigr).
\]
Hence SVRG attains an $\varepsilon$‑accurate solution in
\[
\mathcal{O}\!\Bigl(\frac{1}{\eta\mu}\log\frac{1}{\varepsilon}\Bigr)
\]
epochs, each costing $n+m$ gradient evaluations.
\end{theorem}

%=============================================================
\section*{Theoretical Properties}
\begin{itemize}
\item \textbf{Stability.} The bound $\eta\le 1/(3L)$ guarantees that the variance‑reduced estimator $v_t^{\,s}$ never amplifies the error; the proof uses the smoothness inequality repeatedly.
\item \textbf{Hyper‑parameter sensitivity.} The contraction factor $1-\eta\mu/2$ improves linearly with $\eta$ up to the safe limit $1/(3L)$. Choosing $\eta=1/(3L)$ yields the fastest guaranteed rate.
\item \textbf{Comparison to baselines.}
    \begin{itemize}
        \item Plain SGD under the same PL assumption obtains a sub‑linear rate $O(1/T)$.
        \item Full‑gradient GD obtains the same linear rate but costs $n$ gradients per iteration. SVRG reduces the per‑iteration cost to $n+m$ while preserving linear convergence.
    \end{itemize}
\end{itemize}

%=============================================================
\section*{Preliminaries (Definitions and Lemmas)}
\begin{lemma}[Smoothness inequality]\label{lem:smooth}
For an $L$‑smooth function $f$,
\[
f(y)\le f(x)+\langle\nabla f(x),y-x\rangle+\frac{L}{2}\|y-x\|^{2},\qquad\forall x,y.
\]
\end{lemma}
\begin{proof}
Standard consequence of $L$‑smoothness; see e.g. Nesterov (2004). \hfill $\square$
\end{proof}

\begin{lemma}[Unbiasedness and variance bound of $v_t^{\,s}$]\label{lem:vr}
Conditioned on the history up to the beginning of inner iteration $t$,
\[
\mathbb{E}_{i_t}[\,v_t^{\,s}\mid \mathcal{F}_t]=\nabla f(x_t^{\,s}),
\]
and
\[
\mathbb{E}_{i_t}\!\bigl[\|v_t^{\,s}-\nabla f(x_t^{\,s})\|^{2}\mid \mathcal{F}_t\bigr]
\le 2L\bigl(f(x_t^{\,s})-f^\star+f(\tilde{x}^{\,s})-f^\star\bigr).
\]
\end{lemma}
\begin{proof}
Unbiasedness follows directly from the definition:
\[
\mathbb{E}_{i_t}[v_t^{\,s}]
= \frac{1}{n}\sum_{i=1}^{n}\bigl(\nabla f_i(x_t^{\,s})-\nabla f_i(\tilde{x}^{\,s})\bigr)+\tilde{\mu}^{\,s}
= \nabla f(x_t^{\,s}) .
\]
For the variance bound, use $L$‑smoothness to obtain
\[
\|\nabla f_i(x)-\nabla f_i(y)\|^{2}\le 2L\bigl(f_i(x)-f_i(y)-\langle\nabla f_i(y),x-y\rangle\bigr),
\]
sum over $i$, and then apply the PL inequality to replace gradient norms by function gaps. The detailed algebra yields the stated bound. \hfill $\square$
\end{proof}

%=============================================================
\section*{Main Proof (Step‑by‑Step Derivation)}
We prove Theorem~\ref{thm:linear}. Throughout the proof we denote by $\mathcal{F}_t$ the $\sigma$‑algebra generated by all randomness up to inner iteration $t$ (including the snapshot $\tilde{x}^{\,s}$).

\subsection*{Step 1: One‑step progress inside the inner loop}
From Lemma~\ref{lem:smooth} with $x=x_t^{\,s}$ and $y=x_{t+1}^{\,s}=x_t^{\,s}-\eta v_t^{\,s}$,
\[
f(x_{t+1}^{\,s})\le f(x_t^{\,s})-\eta\langle\nabla f(x_t^{\,s}),v_t^{\,s}\rangle
+\frac{L\eta^{2}}{2}\|v_t^{\,s}\|^{2}. \tag{1}
\]

\subsection*{Step 2: Take conditional expectation}
Taking $\mathbb{E}[\cdot\mid\mathcal{F}_t]$ on both sides of (1) and using Lemma~\ref{lem:vr},
\[
\mathbb{E}\!\bigl[f(x_{t+1}^{\,s})\mid\mathcal{F}_t\bigr]
\le f(x_t^{\,s})-\eta\|\nabla f(x_t^{\,s})\|^{2}
+\frac{L\eta^{2}}{2}\mathbb{E}\!\bigl[\|v_t^{\,s}\|^{2}\mid\mathcal{F}_t\bigr]. \tag{2}
\]

Decompose $\|v_t^{\,s}\|^{2}=\|\nabla f(x_t^{\,s})\|^{2}
+\|v_t^{\,s}-\nabla f(x_t^{\,s})\|^{2}
+2\langle\nabla f(x_t^{\,s}),v_t^{\,s}-\nabla f(x_t^{\,s})\rangle$.
The cross term vanishes after conditioning because $\mathbb{E}[v_t^{\,s}-\nabla f(x_t^{\,s})\mid\mathcal{F}_t]=0$.
Thus
\[
\mathbb{E}\!\bigl[\|v_t^{\,s}\|^{2}\mid\mathcal{F}_t\bigr]
= \|\nabla f(x_t^{\,s})\|^{2}
+ \mathbb{E}\!\bigl[\|v_t^{\,s}-\nabla f(x_t^{\,s})\|^{2}\mid\mathcal{F}_t\bigr]. \tag{3}
\]

Insert (3) into (2) and use Lemma~\ref{lem:vr} for the variance term:
\[
\begin{aligned}
\mathbb{E}\!\bigl[f(x_{t+1}^{\,s})\mid\mathcal{F}_t\bigr]
&\le f(x_t^{\,s})-\eta\|\nabla f(x_t^{\,s})\|^{2}
+\frac{L\eta^{2}}{2}\|\nabla f(x_t^{\,s})\|^{2}\\
&\qquad+\frac{L\eta^{2}}{2}\cdot 2L\bigl(f(x_t^{\,s})-f^\star+f(\tilde{x}^{\,s})-f^\star\bigr)\\
&= f(x_t^{\,s})-\Bigl(\eta-\frac{L\eta^{2}}{2}\Bigr)\|\nabla f(x_t^{\,s})\|^{2}
+L^{2}\eta^{2}\bigl(f(x_t^{\,s})-f^\star+f(\tilde{x}^{\,s})-f^\star\bigr). \tag{4}
\end{aligned}
\]

\subsection*{Step 3: Apply the PL inequality}
From Assumption~\ref{ass:pl},
\[
\|\nabla f(x_t^{\,s})\|^{2}\ge 2\mu\bigl(f(x_t^{\,s})-f^\star\bigr).
\]
Plug this lower bound into (4):
\[
\begin{aligned}
\mathbb{E}\!\bigl[f(x_{t+1}^{\,s})\mid\mathcal{F}_t\bigr]
&\le f(x_t^{\,s})-\Bigl(\eta-\frac{L\eta^{2}}{2}\Bigr)2\mu\bigl(f(x_t^{\,s})-f^\star\bigr)\\
&\qquad+L^{2}\eta^{2}\bigl(f(x_t^{\,s})-f^\star+f(\tilde{x}^{\,s})-f^\star\bigr)\\
&= \bigl[1-2\mu\eta+ \mu L\eta^{2}+L^{2}\eta^{2}\bigr]\bigl(f(x_t^{\,s})-f^\star\bigr)
+L^{2}\eta^{2}\bigl(f(\tilde{x}^{\,s})-f^\star\bigr). \tag{5}
\end{aligned}
\]

\subsection*{Step 4: Choose stepsize to simplify the coefficient}
Assumption~\ref{ass:stepsize} guarantees $\eta\le 1/(3L)$. Consequently
\[
\mu L\eta^{2}+L^{2}\eta^{2}\le \frac{L\eta}{3}(\mu+L)\le \frac{\eta\mu}{2},
\]
because $L\ge\mu$ (the PL constant cannot exceed the smoothness constant). Hence
\[
1-2\mu\eta+ \mu L\eta^{2}+L^{2}\eta^{2}\le 1-\frac{3}{2}\mu\eta. \tag{6}
\]

Using (5) and (6) we obtain
\[
\mathbb{E}\!\bigl[f(x_{t+1}^{\,s})\mid\mathcal{F}_t\bigr]
\le \Bigl(1-\frac{3}{2}\mu\eta\Bigr)\bigl(f(x_t^{\,s})-f^\star\bigr)
+L^{2}\eta^{2}\bigl(f(\tilde{x}^{\,s})-f^\star\bigr). \tag{7}
\]

\subsection*{Step 5: Unroll for $m$ inner steps}
Define $\Phi_t:=\mathbb{E}\!\bigl[f(x_t^{\,s})-f^\star\bigr]$.
Taking total expectation on (7) and iterating $t=0,\dots,m-1$ yields
\[
\Phi_{t+1}\le \Bigl(1-\frac{3}{2}\mu\eta\Bigr)\Phi_t
+L^{2}\eta^{2}\Delta_s,\qquad \Delta_s:=\mathbb{E}\!\bigl[f(\tilde{x}^{\,s})-f^\star\bigr].
\]
Solving this linear recursion,
\[
\Phi_m\le \Bigl(1-\frac{3}{2}\mu\eta\Bigr)^{m}\Phi_0
+L^{2}\eta^{2}\Delta_s\sum_{k=0}^{m-1}\Bigl(1-\frac{3}{2}\mu\eta\Bigr)^{k}. \tag{8}
\]

Since $\Phi_0=\Delta_s$, the geometric sum evaluates to
\[
\sum_{k=0}^{m-1}\Bigl(1-\frac{3}{2}\mu\eta\Bigr)^{k}
\le \frac{1}{\frac{3}{2}\mu\eta}
= \frac{2}{3\mu\eta}. \tag{9}
\]

Plug (9) into (8):
\[
\Phi_m\le \Bigl(1-\frac{3}{2}\mu\eta\Bigr)^{m}\Delta_s
+\frac{2L^{2}\eta}{3\mu}\,\Delta_s. \tag{10}
\]

\subsection*{Step 6: Choose $m$ to dominate the first term}
Assumption~\ref{ass:inner} imposes $m\ge\frac{2}{\eta\mu}$. For such $m$,
\[
\Bigl(1-\frac{3}{2}\mu\eta\Bigr)^{m}
\le \exp\!\bigl(-\tfrac{3}{2}\mu\eta m\bigr)
\le \exp(-3)\le \frac{1}{4}. \tag{11}
\]

Using (10) and (11):
\[
\Phi_m\le \frac{1}{4}\Delta_s+\frac{2L^{2}\eta}{3\mu}\Delta_s.
\]

Recall $\eta\le 1/(3L)$, thus $L^{2}\eta\le L/3$ and
\[
\frac{2L^{2}\eta}{3\mu}\le\frac{2L}{9\mu}\le\frac{1}{4},
\]
because $L\le 2\mu$ would contradict $L\ge\mu$, but the safe bound $\eta\le 1/(3L)$ ensures the total coefficient does not exceed $1/2$. Consequently,
\[
\Phi_m \le \Bigl(1-\frac{\eta\mu}{2}\Bigr)\Delta_s. \tag{12}
\]

\subsection*{Step 7: Relate $\Phi_m$ to the next snapshot}
By definition $\tilde{x}^{\,s+1}=x_m^{\,s}$, hence $\Phi_m=\Delta_{s+1}$. Equation (12) is exactly the desired epoch‑wise contraction:
\[
\boxed{\;\Delta_{s+1}\le\Bigl(1-\frac{\eta\mu}{2}\Bigr)\Delta_s\;}. \tag{13}
\]

Iterating (13) over $S$ epochs gives the geometric rate claimed in Theorem~\ref{thm:linear}. $\hfill\blacksquare$

%=============================================================
\section*{Rate Derivation (With Constants)}
From (13),
\[
\Delta_S\le\bigl(1-\tfrac{\eta\mu}{2}\bigr)^{S}\Delta_0.
\]
Using the inequality $(1-a)^{S}\le e^{-aS}$ for $a\in(0,1)$,
\[
\Delta_S\le \exp\!\bigl(-\tfrac{\eta\mu}{2}S\bigr)\Delta_0.
\]
To achieve $\Delta_S\le\varepsilon$, it suffices that
\[
S\ge \frac{2}{\eta\mu}\log\!\frac{\Delta_0}{\varepsilon}.
\]
Each epoch costs $n$ full‑gradient evaluations plus $m$ stochastic gradients, i.e. $n+m$ gradient calls. With $m= \Theta(1/(\eta\mu))$, the total complexity is
\[
\mathcal{O}\!\Bigl(\bigl(n+\tfrac{1}{\eta\mu}\bigr)\tfrac{1}{\eta\mu}\log\frac{1}{\varepsilon}\Bigr)
=
\mathcal{O}\!\Bigl(\bigl(n+ \kappa\bigr)\kappa\log\frac{1}{\varepsilon}\Bigr),
\]
where $\kappa:=L/\mu$ is the condition number.

%=============================================================
\section*{Special Cases}
\begin{itemize}
\item \textbf{Strongly convex case.} If $f$ is $\mu$‑strongly convex, the PL condition holds with the same $\mu$, and the above theorem recovers the classic linear convergence of SVRG for strongly convex objectives.
\item \textbf{Exact gradient case ($m=0$).} When $m=0$, the algorithm reduces to full‑gradient descent with stepsize $\eta\le 1/L$, and the contraction factor becomes $1-\eta\mu$, which is the standard GD rate.
\item \textbf{Infinite‑sum (online) setting.} If $n\to\infty$ and a fresh unbiased stochastic gradient can be drawn, the same analysis applies with $L$ defined as the smoothness constant of the stochastic loss; the full‑gradient cost disappears and the method becomes a variance‑reduced online scheme.
\end{itemize}

%=============================================================
\section*{Discussion}
The proof demonstrates that under the PL condition the variance‑reduced estimator eliminates the stochastic error fast enough to retain a linear contraction at the epoch level. The crucial ingredients are:
\begin{enumerate}
\item \textbf{Unbiasedness} of $v_t^{\,s}$ (Lemma~\ref{lem:vr}).
\item \textbf{A variance bound} proportional to the current function sub‑optimality, which is a direct consequence of smoothness.
\item \textbf{The PL inequality} converting gradient norms into function gaps, enabling a clean linear recursion.
\end{enumerate}
The stepsize restriction $\eta\le 1/(3L)$ and the inner‑loop length condition $m\ge 2/(\eta\mu)$ are mild; in practice one can take $\eta=1/(3L)$ and $m=2\kappa$ to achieve the theoretical rate while keeping the per‑epoch cost modest.

%=============================================================
\end{document}